% Generated draft: requires mathematical and source review.
\documentclass[11pt]{article}
\usepackage[margin=25mm]{geometry}
\usepackage{fontspec}
\setmainfont{Latin Modern Roman}
\usepackage{amsmath,amssymb,mathtools,unicode-math}
\setmathfont{Latin Modern Math}
\usepackage{enumitem,needspace,xurl,hyperref}
\setlength{\parindent}{0pt}
\setlength{\parskip}{5pt}
\setlength{\emergencystretch}{3em}
\newcommand{\N}{\mathbb{N}}
\newcommand{\Z}{\mathbb{Z}}
\newcommand{\Q}{\mathbb{Q}}
\newcommand{\R}{\mathbb{R}}
\newcommand{\C}{\mathbb{C}}
\newcommand{\F}{\mathbb{F}}
\newcommand{\A}{\mathbb{A}}
\newcommand{\PP}{\mathbb{P}}
\newcommand{\E}{\mathbb{E}}
\newcommand{\eps}{\varepsilon}
\newcommand{\dd}{\,\mathrm{d}}
\begin{document}
\section*{UnsolvedMath ID: 30003156}
\textbf{Source problem code:} catalogue entry
\subsection{Definitions and mathematical statement}
A symplectic form on a smooth real four-dimensional manifold \(M\) is a smooth real differential two-form \(\omega\) satisfying \(\dd\omega=0\) and nondegeneracy at every point. It determines the orientation for which \(\omega^2:=\omega\wedge\omega\) is positive. Write \([\omega]\in H^2_{\mathrm{dR}}(M;\R)\) for its de Rham cohomology class. For a closed oriented four-manifold, \(b_2^+(M)\) denotes the positive index of the intersection form on \(H^2(M;\R)\).

\textbf{Question as supplied.} For every compact symplectic four-manifold denoted \((M,J,\omega)\) with \(b_2^+(M)=1\), and every smooth real-valued function \(F\in C^\infty(M;\R)\) satisfying
\[
\int_M e^F\omega^2=\int_M\omega^2,
\]
does there exist a symplectic form \(\widetilde{\omega}\) on \(M\) such that
\[
[\widetilde{\omega}]=[\omega]\quad\text{in }H^2_{\mathrm{dR}}(M;\R),
\qquad
\widetilde{\omega}\wedge\widetilde{\omega}
=e^F\omega\wedge\omega?
\]
Here \(\widetilde{\omega}\in[\omega]\) in the source is interpreted as equality of de Rham cohomology classes.

\textbf{Unresolved qualifications.} The source includes \(J\) without defining it or imposing any relation between \(J\), \(\omega\), and \(\widetilde{\omega}\). If \(J\) is intended to be an almost complex structure, its defining condition is \(J^2=-\mathrm{id}_{TM}\). Compatibility of a form \(\eta\) with such a \(J\) would require \(\eta(Ju,Jv)=\eta(u,v)\) and \(\eta(v,Jv)>0\) for nonzero \(v\); taming would require the latter positivity condition. Neither requirement is explicitly present in the supplied question, and neither is added here. Consequently the displayed question does not specify the intended role of \(J\) in the problem named by the title. The source also does not explicitly say that \(M\) is connected or has empty boundary. The definition of \(b_2^+\) above uses the closed-manifold interpretation; a boundary interpretation and any boundary conditions remain unspecified. If disconnected manifolds are allowed, the source supplies only the displayed total-volume normalization, with no separate normalization on each component.
\subsection{Short English statement}
The supplied question asks whether every normalized density \(e^F\omega^2\) can be realized as the square of a symplectic form cohomologous to \(\omega\) on a compact four-manifold with \(b_2^+=1\). The role of \(J\), any compatibility or taming requirements, and the intended connectedness and boundary assumptions require clarification before this can serve as a precise formulation of the named Donaldson problem.
\subsection{Sources}
\begin{itemize}
\item Supplied local UnsolvedMath registry record, permanent ID \texttt{30003156}, source problem code \texttt{catalogue entry}, title \emph{Donaldson Symplectic Calabi--Yau Equation}. The project identifies \texttt{data/unsolvedmath/problems.json} and the supplementary \texttt{catalog/problems.json} as local registries; their contents were not independently inspected for this draft. The record gives proposed year 2016 and names no proposer.
\item Supplied report citation: \emph{Geometrie (2017)}, Oberwolfach Reports, pp.\ 1577--1634, DOI 10.4171/owr/2016/28. \url{https://doi.org/10.4171/owr/2016/28}. The background identifies this report as documenting workshop open-problem discussions.
\item Supplied literature source: Weinkove, \emph{The Calabi--Yau equation on almost Kahler manifolds}. \url{https://arxiv.org/abs/math/0604408}.
\item Supplied literature source: Fino--Li--Salamon--Vezzoni. \url{https://doi.org/10.1090/S0002-9947-2012-05692-3}. No paper title is supplied.
\item Supplied literature source: Oberwolfach Report 28. \url{https://doi.org/10.4171/owr/2016/28}.
\item Literature assessment dated \textbf{2026-08-21}: the supplied curation labels the problem open, reports existence results of Weinkove under small-Nijenhuis assumptions and a reduction of the general problem to a priori estimates, and says that later work treats special cases while the general Donaldson symplectic Calabi--Yau equation remains unsolved. This is supplied dated curation derived from the OWR material, not an independently verified assessment. The source explicitly requests independent verification of time-sensitive status and cited claims.
\end{itemize}
\end{document}
